% RQ1 Abstract: Participation Dynamics

Low voter turnout is a persistent challenge in Decentralized Autonomous Organizations (DAOs), threatening governance legitimacy and operational effectiveness. We present a systematic investigation of how participation incentives and quorum calibration affect voter turnout and retention using multi-agent simulations.

Through \experimentruns{} Monte Carlo simulation runs across two experiment sets, we examine the relationship between participation target rates, quorum thresholds, and long-term voter engagement. Our simulations model heterogeneous agent populations with varying participation propensities, token distributions, and fatigue dynamics.

Key findings include: (1) participation targets set above natural engagement rates create unsustainable quorum requirements leading to governance gridlock; (2) there exists a non-linear relationship between quorum thresholds and effective governance, with an optimal ``sweet spot'' balancing legitimacy and operability; (3) voter fatigue accumulates faster under high-frequency voting, suggesting benefits to proposal batching or conviction-based mechanisms.

We provide actionable calibration guidelines for DAO practitioners and release our simulation framework as open-source software to enable reproducible governance research.
